% GENERATED FILE. Edit canonical lessons, then run npm run generate:books.
% canonical-source-hash: fnv1a64:43b70afd9fbc93ff
% canonical-lessons: TE-C191-rabbaru, TE-C191-balbu, TE-C191-kattelu, TE-C191-vala, TE-C191-trasu

\chapter{Rubber, Light bulb, Firewood, Net, Weighing scale}
\label{ch:te-rubber-light-bulb-firewood-net-weighing-scale}
\hlchaptermodality{191}

\begin{chapteropening}
\textbf{By the end of this chapter:} \emph{I can say rubber, a light bulb, firewood, a net and a weighing scale.}

Complete the last lesson of chapter 191: I can say rubber, a light bulb, firewood, a net and a weighing scale.
\end{chapteropening}

\section[\texorpdfstring{\te{రబ్బరు} (rabbaru)}{రబ్బరు (rabbaru)}]{\te{రబ్బరు} (rabbaru) — rubber; an eraser}

\label{lesson:TE-C191-rabbaru}



\begin{quote}
Before the new one: say the Telugu for a swing, then the Telugu for a carpet.
\end{quote}

\subsection*{You'll want to know: \te{రబ్బరు}}
\textbf{\te{రబ్బరు}} — \emph{rabbaru} — \textquotedblleft{}rubber; an eraser\textquotedblright{}.

\begin{grammarlens}[title={What you've built}]
One more word for this chapter.
\end{grammarlens}

\subsection*{Your turn}
\begin{itemize}
\raggedright
  \item \emph{Say it:} \emph{rabbaru}
  \item \emph{Say it:} \emph{rabbaru}, once more
  \item \emph{Say it:} say \emph{rabbaru}
  \item \emph{Recall:} say the Telugu for a swing, then the Telugu for a carpet, then say \emph{rabbaru} again
\end{itemize}

\begin{tcolorbox}[breakable,colback=teal!4,colframe=teal!35!black,title={Before you move on}]
What does \te{రబ్బరు} mean? (\textquotedblleft{}Rubber; an eraser\textquotedblright{}.) Say it once more.
\end{tcolorbox}
\section[\texorpdfstring{\te{బల్బు} (balbu)}{బల్బు (balbu)}]{\te{బల్బు} (balbu) — a light bulb}

\label{lesson:TE-C191-balbu}



\begin{quote}
Before the new one: say the Telugu for a carpet, then the Telugu for rubber.
\end{quote}

\subsection*{You'll want to know: \te{బల్బు}}
\textbf{\te{బల్బు}} — \emph{balbu} — \textquotedblleft{}a light bulb\textquotedblright{}.

\begin{grammarlens}[title={What you've built}]
One more word for this chapter.
\end{grammarlens}

\subsection*{Your turn}
\begin{itemize}
\raggedright
  \item \emph{Say it:} \emph{balbu}
  \item \emph{Say it:} \emph{balbu}, once more
  \item \emph{Say it:} say \emph{balbu}
  \item \emph{Recall:} say the Telugu for a carpet, then the Telugu for rubber, then say \emph{balbu} again
\end{itemize}

\begin{tcolorbox}[breakable,colback=teal!4,colframe=teal!35!black,title={Before you move on}]
What does \te{బల్బు} mean? (\textquotedblleft{}A light bulb\textquotedblright{}.) Say it once more.
\end{tcolorbox}
\section[\texorpdfstring{\te{కట్టెలు} (kaṭṭelu)}{కట్టెలు (kaṭṭelu)}]{\te{కట్టెలు} (kaṭṭelu) — firewood}

\label{lesson:TE-C191-kattelu}



\begin{quote}
Before the new one: say the Telugu for rubber, then the Telugu for a light bulb.
\end{quote}

\subsection*{You'll want to know: \te{కట్టెలు}}
\textbf{\te{కట్టెలు}} — \emph{kaṭṭelu} — \textquotedblleft{}firewood\textquotedblright{}.

\begin{grammarlens}[title={What you've built}]
One more word for this chapter.
\end{grammarlens}

\subsection*{Your turn}
\begin{itemize}
\raggedright
  \item \emph{Say it:} \emph{kaṭṭelu}
  \item \emph{Say it:} \emph{kaṭṭelu}, once more
  \item \emph{Say it:} say \emph{kaṭṭelu}
  \item \emph{Recall:} say the Telugu for rubber, then the Telugu for a light bulb, then say \emph{kaṭṭelu} again
\end{itemize}

\begin{tcolorbox}[breakable,colback=teal!4,colframe=teal!35!black,title={Before you move on}]
What does \te{కట్టెలు} mean? (\textquotedblleft{}Firewood\textquotedblright{}.) Say it once more.
\end{tcolorbox}
\section[\texorpdfstring{\te{వల} (vala)}{వల (vala)}]{\te{వల} (vala) — a net}

\label{lesson:TE-C191-vala}



\begin{quote}
Before the new one: say the Telugu for a light bulb, then the Telugu for firewood.
\end{quote}

\subsection*{You'll want to know: \te{వల}}
\textbf{\te{వల}} — \emph{vala} — \textquotedblleft{}a net\textquotedblright{}.

\begin{grammarlens}[title={What you've built}]
One more word for this chapter.
\end{grammarlens}

\subsection*{Your turn}
\begin{itemize}
\raggedright
  \item \emph{Say it:} \emph{vala}
  \item \emph{Say it:} \emph{vala}, once more
  \item \emph{Say it:} say \emph{vala}
  \item \emph{Recall:} say the Telugu for a light bulb, then the Telugu for firewood, then say \emph{vala} again
\end{itemize}

\begin{tcolorbox}[breakable,colback=teal!4,colframe=teal!35!black,title={Before you move on}]
What does \te{వల} mean? (\textquotedblleft{}A net\textquotedblright{}.) Say it once more.
\end{tcolorbox}
\section[\texorpdfstring{\te{త్రాసు} (trāsu)}{త్రాసు (trāsu)}]{\te{త్రాసు} (trāsu) — a weighing scale}

\label{lesson:TE-C191-trasu}



\begin{quote}
Before the new one: say the Telugu for firewood, then the Telugu for a net.
\end{quote}

\subsection*{You'll want to know: \te{త్రాసు}}
\textbf{\te{త్రాసు}} — \emph{trāsu} — \textquotedblleft{}a weighing scale\textquotedblright{}.

\begin{grammarlens}[title={What you've built}]
One more word for this chapter.
\end{grammarlens}

\subsection*{Your turn}
\begin{itemize}
\raggedright
  \item \emph{Say it:} \emph{trāsu}
  \item \emph{Say it:} \emph{trāsu}, once more
  \item \emph{Say it:} say \emph{trāsu}
  \item \emph{Recall:} say the Telugu for firewood, then the Telugu for a net, then say \emph{trāsu} again
\end{itemize}

\begin{tcolorbox}[breakable,colback=teal!4,colframe=teal!35!black,title={Before you move on}]
What does \te{త్రాసు} mean? (\textquotedblleft{}A weighing scale\textquotedblright{}.) Say it once more.
\end{tcolorbox}
